\section{Leading equivalents and the critical window}\label{ssm:sec:results}
Let
\[
 \phi(x)=\frac{e^{-x^2/2}}{\sqrt{2\pi}},\qquad
 \Phi(x)=\int_{-\infty}^x\phi(t)\,\dd t,
\]
and let $\ell$ be the unique positive solution
\begin{equation}\label{ssm:eq:ell}
 \phi(\ell)=\ell\Phi(\ell).
\end{equation}
To see existence and uniqueness directly, $h(x)=\phi(x)-x\Phi(x)$ has $h(0)>0$ and $h'(x)=-\Phi(x)-2x\phi(x)<0$ for $x\ge0$, while $h(1)<0$. In fact $0<\ell<1/\sqrt2$: at $x=1/\sqrt2$, the bound $\Phi(x)>1/2+x\phi(x)$ gives $x\Phi(x)>\phi(x)$ because $x>\phi(x)$. Define
\begin{equation}\label{ssm:eq:constants}
 \begin{split}
 q&=\Phi(\ell)e^{-\ell^2/2},\qquad
 Z_0=\frac{\exp(5\ell^2/4-\ell^4/2)}{\sqrt{1+2\ell^2}},\\
 J&=-\frac{4\ell^2}{1+2\ell^2},\qquad
 v=\frac{1-2\ell^2}{4}\in(0,1/4).
 \end{split}
\end{equation}

\begin{theorem}[Mixed degree constraints]\label{ssm:thm:transfer}
Fix finite constants $C,K\ge0$. Let $M\to\infty$, let $k$ be integral with $|k-M/2|\le K$, and let $0\le f\le C\sqrt M$ be integral. Choose the finite saddle $r_M=1/2-L_M/(2\sqrt M)$ described in Section~\ref{ssm:sec:saddle}, and put
\[
 S_M=\PP\{\Bin(M-1,r_M)\le k\},\qquad
 \theta_M=k-M/2+1.
\]
Then, uniformly over these $k$ and $f$,
\begin{equation}\label{ssm:eq:criticalexact}
 \frac{U_f(M,k)}{B(M,k)}
 =(1+o(1))S_M^{-f}
 \exp\left\{\frac{\ell^2}{2(1+2\ell^2)}\frac{f^2}{M}\right\}.
\end{equation}
Equivalently, with no implicit saddle mass,
\begin{equation}\label{ssm:eq:criticalexplicit}
 \frac{U_f(M,k)}{B(M,k)}
 =(1+o(1))\Phi(\ell)^{-f}
 \exp\left\{-\frac{2\theta_M\ell}{1+2\ell^2}\frac f{\sqrt M}
       +\frac{\ell^2}{2(1+2\ell^2)}\frac{f^2}{M}\right\}.
\end{equation}
In particular, $U_f/B\sim\Phi(\ell)^{-f}$ if $f=o(\sqrt M)$, including every fixed $f$.
\end{theorem}

The equivalence of \eqref{ssm:eq:criticalexact} and \eqref{ssm:eq:criticalexplicit} requires the source-based estimate
\begin{equation}\label{ssm:eq:Srateintro}
 \log S_M=\log\Phi(\ell)+\frac{2\theta_M\ell}{1+2\ell^2}M^{-1/2}
 +O(M^{-1}),
\end{equation}
proved in Section~\ref{ssm:sec:explicit}. The mere convergence $S_M\to\Phi(\ell)$ would not suffice when $f$ grows like $\sqrt M$. If $f/\sqrt M\to\tau$, the quadratic term has the nonvanishing limit $\ell^2\tau^2/[2(1+2\ell^2)]$. The linear term remembers the bounded offset of the cap. Both matter on this scale.

\begin{theorem}[Symmetric sign matrices]\label{ssm:thm:matrix}
With the constants in \eqref{ssm:eq:constants},
\begin{equation}\label{ssm:eq:matrix}
 A_n\sim 2^{n(n+1)/2}q^n Z_0e^{-\ell^2/2}
 \begin{cases}
 \exp(\ell\sqrt n+J/4),&n\text{ even},\\
 1,&n\text{ odd}.
 \end{cases}
\end{equation}
Each equivalence is a relative $1+o(1)$ statement along its parity subsequence.
\end{theorem}

Useful decimal values, displayed for orientation, are
\[
\begin{array}{rlrl}
 \ell&\approx0.50605446898918076324,&\Phi(\ell)&\approx0.69359079375820639821,\\
 q&\approx0.61023039919762778828,&Z_0&\approx1.08387834752412285047,\\
 J&\approx-0.67740809782986162079,&\Phi(\ell)^{-1}&\approx1.44177230868582047454.
\end{array}
\]
Ordinary high-precision decimal evaluations and any exact interval certificates in the computational package are distinguished there. The decimal display itself is not used to establish the theorem.

The diagonal convention is essential. It is responsible for the extra free vertex in \eqref{ssm:eq:bridge}; deleting or fixing the diagonal produces a different count. The parity factor is also essential: omitting it from even orders would not yield a relative equivalent. Neither theorem asserts a rate for its $o(1)$, higher asymptotic count coefficients, or a numerically certified starting index.
